% =============================================================
% BODY (EN)
% =============================================================

\section{Activity Solution}
\vspace{2mm}

The selected paper for this activity is \textit{Sentiment Analysis in Social Networks and Online Search Engines: A Financial Predictive Approach}. The paper analyzes how sentiment analysis, applied to social networks, online search engines, financial news and other digital sources, can complement traditional financial data in order to better understand and predict market behavior.

\vspace{2mm}

The following sections answer the four questions proposed in the activity, based on the reading and analysis of the selected paper.


\subsection{What Was The Problem?}
\vspace{2mm}

The main problem identified in the paper is the difficulty of accurately predicting variations in financial markets. These variations are not explained only by historical data or traditional economic indicators, but also by several external factors, such as global economic changes, news, investor behavior and social perceptions expressed through digital platforms.

\vspace{2mm}

The paper explains that many financial operators assume that historical performance can be used to predict future returns. However, the work expands this view by arguing that user sentiment also plays an important role. Therefore, the problem is not only statistical or financial, but also informational: traditional models may be incomplete if they do not include non-structured data from social networks, online searches, financial news or forums.

\vspace{2mm}

The paper also highlights that consumers and investors are increasingly difficult to predict. For this reason, it becomes necessary to include new sources of information capable of reflecting opinions, emotions and trends in real time. From a software design perspective, the problem can be interpreted as the need to build solutions that integrate heterogeneous data sources, process natural language and generate useful information for financial decision-making.


\subsection{What Solution Does The Paper Propose?}
\vspace{2mm}

The proposed solution is not a single implemented software tool, but an approach based on a systematic mapping of the literature. The goal of the paper is to review the current state of sentiment analysis techniques applied to social networks and online search engines, with a specific focus on financial prediction.

\vspace{2mm}

The paper proposes the integration of sentiment analysis with financial predictive models. This integration allows traditional financial data to be complemented with non-structured information obtained from digital platforms. Sentiment analysis classifies texts according to their polarity, for example positive, negative or neutral, and this information can then be transformed into an additional signal to interpret market trends.

\vspace{2mm}

The solution is based on the idea that social networks, web news, Google trends and forum discussions may reflect collective perceptions about companies, assets or economic contexts. By combining these signals with predictive models, the approach seeks to improve decision-making, anticipate market movements and support risk management more quickly.

\vspace{2mm}

The paper also indicates that predictive accuracy can improve when multiple data sources are combined, such as news, social media and other online platforms. Therefore, the general proposal is to design an analysis approach that does not depend on a single source, but integrates traditional financial information with textual data from the digital environment.


\subsection{What Design Decisions Appear?}
\vspace{2mm}

Although the paper does not present the detailed design of a specific application, it does reveal several relevant design decisions for a system based on sentiment analysis applied to financial prediction.

\vspace{2mm}

First, the paper adopts a \textbf{systematic mapping of the literature} as its research strategy. This decision implies defining research questions, selecting sources, building search strings and applying inclusion and exclusion criteria. From a design point of view, this provides structure, traceability and justification for the results obtained.

\vspace{2mm}

Second, the paper decides to work with \textbf{multiple digital sources}. For the literature review, it mentions sources such as SEDICI, IEEE Xplore Digital Library, ACM Digital Library and Science Direct. For a predictive solution, it also considers social networks, web news, Google trends and financial forum discussions. This decision aims to enrich the data set and avoid relying only on traditional financial indicators.

\vspace{2mm}

Third, there is a design decision related to the use of \textbf{natural language processing}. The system must transform opinions, posts, news and comments into structured information that can be used by predictive models. This implies defining mechanisms for extraction, cleaning, classification and representation of textual data.

\vspace{2mm}

Another important decision is the \textbf{domain delimitation}. The paper does not study sentiment analysis in a general way, but focuses specifically on social networks, online search engines and financial prediction. This scope helps organize the research questions and makes the analysis more precise.

\vspace{2mm}

Finally, the paper considers both \textbf{benefits and limitations}. Sentiment analysis is not presented as an automatic or perfect solution. The paper recognizes challenges related to model complexity, language ambiguity, market volatility, transparency, privacy and the need for financial analysts to understand and use these techniques properly.


\subsection{Which Quality Attributes Does It Try To Improve?}
\vspace{2mm}

The most evident quality attribute addressed by the proposed approach is \textbf{accuracy}. The paper suggests that incorporating information from social networks, search engines and financial news can complement traditional financial data, allowing more precise predictions about market trends.

\vspace{2mm}

Another important attribute is \textbf{timeliness}. By using digital sources that can provide information in real time, the approach may detect opinion changes, emerging trends or public reactions before they are fully reflected in traditional financial indicators. This can support faster decisions by investors, analysts and companies.

\vspace{2mm}

The paper also relates to \textbf{adaptability}. It explains that sentiment analysis can help personalize financial strategies according to investor profiles, risk preferences and behaviors. Therefore, a solution based on this approach should be able to adapt its recommendations to different users and decision contexts.

\vspace{2mm}

\textbf{Robustness} is also relevant. The paper recognizes that social media and news language can be ambiguous, and that financial markets are volatile and complex. For this reason, a system of this type must be able to distinguish relevant information from noise, handle multiple data sources and avoid weak conclusions based on unstable or subjective information.

\vspace{2mm}

Another quality attribute is \textbf{security and privacy}. In its conclusions, the paper mentions the need to guarantee the privacy and security of financial data. This is especially important because a financial predictive system may work with sensitive information from users, investors or institutions.

\vspace{2mm}

Finally, the paper points indirectly to \textbf{usability} and \textbf{understandability}. It warns that the complexity of predictive models must be balanced with transparency and user comprehension. In other words, a system should not only predict market behavior, but also present its results clearly enough to support responsible decision-making.


\section{Conclusion}
\vspace{2mm}

In summary, the paper addresses a current and relevant problem: the difficulty of anticipating financial behavior in highly variable contexts. The proposed approach consists of studying and integrating sentiment analysis techniques with predictive models, using non-structured information from social networks, search engines, financial news and other digital sources.

\vspace{2mm}

From a software design perspective, the paper makes it possible to identify decisions related to source selection, natural language processing, data integration, domain delimitation and risk consideration. The main quality attributes involved are accuracy, timeliness, adaptability, robustness, security, privacy and clarity in the presentation of results.